\begin{textsample}{POS Dim 4 – vem\_ed\_15 – Score 2.00 – vem\_ed\_15/t064.txt}  \label{ex:f4_pos_119}
IVES: ABERTURA Em 8 de dezembro, das 9 às 16h, equipes de Coordenação de Projetos da Cesu e a equipe dos Projetos Colaborativos Internacionais (PCIs / Cesu) organizaram o IVES Cesu 2022 - International Virtual Exchange Symposium (Simpósio Internacional de Intercâmbios Virtuais).
Participaram, como ouvintes, 176 pessoas de 11 países, incluindo o Brasil.
O evento está disponível no Canal do YouTube da Cesu: youtube.com/playlist? list = PLF0wjfo4cP3Hvg1eNwxPiMHFje3\_J9L7P Osvaldo Succi Junior, coordenador dos PCIs / Cesu, foi responsável pela mediação do evento, que contou com cerimonial de Renata Cardias (Cesu) e Regiane Camargo (PCIs / Cesu).
Na abertura, Rafael Ferreira Alves, coordenador técnico da Cesu, enfatizou a importância dos PCIs para o processo de internacionalização das Fatecs.
Desde 2013, participaram dos PCIs 7.400 alunos em mais de 50 Fatecs, com colaborações em 15 países diferentes.
Somente no ano de 2022, foram 125PCIs, com 2.633 alunos.
O Centro Paula Souza é reconhecido internacionalmente pela realização de Intercâmbios Virtuais e, em outubro de 2022, recebeu o selo BraVE (acrônimo para Brazilian Virtual Exchange).
Trata-se de um programa da Associação Brasileira de Educação Internacional - FAUBAI que incentiva a qualificação das Instituições de Ensino Superior associadas para a realização de Intercâmbios Virtuais.
Marta Iglesis Farrero, assessora de relações internacionais da ARInter - CPS, destacou que a internacionalização é uma ferramenta estratégica para alcançar a excelência de alunos, professores e servidores \textbf{administrativos}.
André Luiz Braun Galvão, diretor do Departamento Acadêmico Pedagógico da Cesu, reconheceu a importância dos PCIs para o desenvolvimento profissional dos estudantes.
Mariane Teixeira, coordenadora de Línguas e Projetos Internacionais da Cesu, ressaltou a importância da parceria entre ensino de línguas nos cursos das Fatecs e a realização dos Projetos Colaborativos Internacionais, por professores de línguas e de disciplinas técnicas, para a internacionalização do currículo.
Os Intercâmbios Virtuais ou, como são conhecidos mundialmente, Collaborative Online International Learning (COIL), propõem o trabalho colaborativo de professores de diferentes países na \textbf{criação} de projetos e tarefas a serem realizadas conjuntamente entre alunos de dois ou mais países, compondo o plano de ensino de suas disciplinas.
No CPS, são chamados de Projetos Colaborativos Internacionais e estão vinculados à Coordenação de Línguas e Projetos Internacionais da Unidade de Ensino Superior de Graduação (PCIs / Cesu).

% matched lemmas: administrativo, criação
\end{textsample}
